% Intended LaTeX compiler: pdflatex
\documentclass[final,12pt]{elsarticle}
\usepackage[utf8]{inputenc}
\usepackage[T1]{fontenc}
\usepackage{graphicx}
\usepackage{longtable}
\usepackage{wrapfig}
\usepackage{rotating}
\usepackage[normalem]{ulem}
\usepackage{amsmath}
\usepackage{amssymb}
\usepackage{capt-of}
\usepackage{hyperref}
\usepackage[utf8]{inputenc}
\usepackage[T1]{fontenc}
\usepackage{lineno}
\linenumbers
\modulolinenumbers[1]
\usepackage{color}
\usepackage{hyperref,xspace}
\usepackage[tight,footnotesize]{subfigure}
\usepackage{tabularx}
\date{}
\title{Using Mandelbrot Set as a Pedagogic Tool in the Teaching of Parallel Computing Concepts}
\hypersetup{
 pdfauthor={Rayan Raddatz de Matos, Francisco Pegoraro Etcheverria, Kenichi Brumati, Lucas Mello Schnorr},
 pdftitle={Using Mandelbrot Set as a Pedagogic Tool in the Teaching of Parallel Computing Concepts},
 pdfkeywords={},
 pdfsubject={},
 pdfcreator={Emacs 30.1 (Org mode 9.7.11)}, 
 pdflang={Portuges}}
\begin{document}

\journal{Information Sciences}
%% use the tnoteref command within \title for footnotes;
%% use the tnotetext command for theassociated footnote;
%% use the fnref command within \author or \address for footnotes;
%% use the fntext command for theassociated footnote;
%% use the corref command within \author for corresponding author footnotes;
%% use the cortext command for theassociated footnote;
%% use the ead command for the email address,
%% and the form \ead[url] for the home page:
%% \title{Title\tnoteref{label1}}
%% \tnotetext[label1]{}
%% \author{Name\corref{cor1}\fnref{label2}}
%% \ead{email address}
%% \ead[url]{home page}
%% \fntext[label2]{}
%% \cortext[cor1]{}
%% \address{Address\fnref{label3}}
%% \fntext[label3]{}

\author[UFRGS]{Rayan Raddatz de Matos}
\ead{rayan.raddatz@inf.ufrgs.br}


\author[UFRGS]{Francisco Pegoraro Etcheverria}
\ead{fpetcheverria@inf.ufrgs.br}


\author[UFRGS]{Kenichi Brumati}
\ead{kbrumati@inf.ufrgs.br}



\author[UFRGS]{Lucas Mello Schnorr}
\ead{schnorr@inf.ufrgs.br}

\address[UFRGS]{
Institute of Informatics, Federal University of Rio Grande do Sul -- UFRGS\\
91501-970, Porto Alegre, RS -- Brazil\\
}
\begin{abstract}

Put the abstract here.

\end{abstract}
%\begin{graphicalabstract}

%\end{graphicalabstract}
%\begin{highlights}

%\end{highlights}
\begin{keyword}

%% keywords here, in the form: keyword \sep keyword

%% PACS codes here, in the form: \PACS code \sep code

%% MSC codes here, in the form: \MSC code \sep code
%% or \MSC[2008] code \sep code (2000 is the default)

\end{keyword}
\maketitle
\section{Introduction}
\label{sec:orgb89c5ab}

\textbf{{[}Context: HPC, CS and Education]}


\textbf{{[}Talk about tools and the problem]}


\textbf{{[}Introduce our tool and explain why use mandelbrot]}


\textbf{{[}Talk about contributions]}


\textbf{{[}Other Sections]}

\textbf{{[}Open science, url, etc]}
\section{Related Work}
\label{sec:orge600160}

{[}Other tools for HPC/PP teaching]

{[}MandelBrot as a default benchmark]

{[}Visualization]

{[}Differences of our approach]
\section{Material and Methods}
\label{sec:org401255e}

\subsection{Application}
\label{sec:orgc458059}

\subsubsection{Arquitecture}
\label{sec:orga15d8ab}

\subsubsection{Client}
\label{sec:orgb8a5364}
\subsection{Visualizing parallel concepts}
\label{sec:org02d7b4c}

{[}Granularity]

{[}Load Balance]

{[}Speedup and Efficiency]

{[}Scalability]
\section{Pedagogic Deployment and Experience}
\label{sec:orgb50066f}

\section{Related Work and Motivation}
\label{sec:orgb740105}
\label{sec.relatedwork}
\section{Your Great Contribution to Computer Science}
\label{sec:orge1642af}
\label{sec.proposal}
\section{Your Super Results}
\label{sec:orge73f018}
\section{Conclusion}
\label{sec:org9dab3fb}
\section*{Acknowledgements}

Who paid for this?
\bibliographystyle{elsearticle-harv}
\bibliography{refs}
\end{document}
